\chapter[反馈线性化]{第 3 讲\quad 部分反馈线性化：把"力"变成"加速度"}
\label{ch:fbl}

\begin{quote}\small
\textbf{本讲目标}\quad 换一个坐标看同一台机器：把"我能施加多大力"换成"我能让小车产生多大加速度"。
小车通道是完全驱动的，所以这个换元是\textbf{精确}的、全域可逆的。换完之后，
后面所有控制器都只需要对付摆杆那一个非线性方程。\\[2pt]
\textbf{前置}\quad 第 1 讲的状态方程；第 2 讲的"局部最优"概念。\\
\textbf{视频}\quad 约 3 分钟。\quad\textbf{代码}\quad \texttt{build/controllers.py: cart\_inverse\_dynamics}。
\end{quote}

\section{为什么值得换坐标}

回顾第 1 讲的两个方程：
\begin{align}
  (m_c+m_p\sin^2\theta)\,\ddot p &= u - b_c\dot p + m_p l\sin\theta\,\dot\theta^2
     + b_p\cos\theta\,\dot\theta - m_p g\sin\theta\cos\theta,\\
  \ddot\theta &= \frac{g\sin\theta - \ddot p\cos\theta - b_p\dot\theta/(m_p l)}{l}.
\end{align}
对控制器设计者来说，这两个式子很讨厌：摩擦、科氏项、重力耦合全混在一起，
每设计一个控制律都要把这些项重新算一遍。

但如果把视角从"力 $u$"换成"小车加速度 $\ddot p$"，情况完全不同：
上式左端的分母恒正（第 1 讲命题 1），于是\textbf{对任意给定的 $\ddot p$，
都有一个唯一的力与之对应}。把这件事写成映射：

\begin{equation}
  u = \underbrace{\big(m_c+m_p\sin^2\theta\big)}_{\text{等效质量}}\,\ddot p
      + \underbrace{b_c\dot p - m_p l\sin\theta\,\dot\theta^2
        - b_p\cos\theta\,\dot\theta + m_p g\sin\theta\cos\theta}_{\text{需要抵消的"干扰"}} .
  \label{eq:inv}
\end{equation}

\keypoint{式 \eqref{eq:inv} 是\textbf{精确}的，不是近似。
它把系统拆成两层：内层负责"抵消已知的非线性"，外层面对一个干净的输入 $v=\ddot p$。
这就是\textbf{部分反馈线性化}（partial feedback linearization）——
"部分"是因为只线性化了小车通道，摆杆通道依旧是非线性的。}

\section{相对阶：为什么只挑小车通道}

对输出 $y=p$，求导两次就出现 $u$（因为 $\ddot p$ 显含 $u$），所以相对阶是 2。
对输出 $y=\theta$，有
\begin{equation}
  \ddot\theta = \frac{g\sin\theta - \ddot p\cos\theta}{l} - \frac{b_p\dot\theta}{m_p l^2},
\end{equation}
其中 $\ddot p$ 显含 $u$，但系数是 $\cos\theta$。因此\textbf{只要 $\cos\theta\ne0$，
$\theta$ 的相对阶也是 2}。

两个输出都是 2 阶，系统总阶是 4，输入只有 1 个。解耦矩阵
\begin{equation}
  \Delta(x) = \begin{bmatrix}
    \dfrac{1}{m_c+m_p\sin^2\theta}\\[8pt]
    \dfrac{-\cos\theta}{l\,(m_c+m_p\sin^2\theta)}
  \end{bmatrix}\in\mathbb{R}^{2\times1}
\end{equation}
是一个\textbf{列向量}，秩恒为 1。所以不可能同时把 $\ddot p$ 和 $\ddot\theta$ 指定成任意值——
这就是第 1 讲说过的欠驱动。工程上的处理是：\textbf{只线性化能线性化的那个通道，
把另一个通道交给"能量"或"滑模"去管}。

\section{反解之后系统长什么样}

把 \eqref{eq:inv} 代入原方程，得到
\begin{empheq}[box=\fbox]{align}
  \ddot p &= v \qquad(\text{单位增益，精确线性化})\\
  \ddot\theta &= \frac{g\sin\theta - v\cos\theta}{l}
     - \frac{b_p\dot\theta}{m_p l^2}
  \label{eq:after}
\end{empheq}
其中 $v$ 是新的控制输入（加速度指令）。

\begin{itemize}[leftmargin=2em,itemsep=2pt]
  \item \textbf{小车通道变成二重积分器}：$p^{(2)}=v$，这是控制理论里最好对付的对象。
  \item \textbf{摆杆通道保留全部非线性}，但形式变干净了：只剩 $g\sin\theta$ 和 $v\cos\theta$，
        以及一个很小的粘性项。后面第 4、5 讲都在这一个方程上做文章。
\end{itemize}

\section{精度验证}

反解是解析的，所以精度只受两个因素影响：浮点误差与\textbf{限幅}。
数值实验（$200$ 组随机状态与随机指令 $v$，用 $10^{-4}$\,s 步长估计实际 $\ddot p$）：

\begin{center}
\small
\begin{tabular}{@{}lcc@{}}
\toprule
指标 & 数值 & 说明 \\
\midrule
最大误差 $|\ddot p - v|$ & $9.2\times10^{-4}$\,m/s$^2$ & 出现在指令被 $\pm12$\,N 限幅的样本 \\
平均误差 & $2.2\times10^{-4}$\,m/s$^2$ & 相当于 $0.02\%$ 的相对误差 \\
\bottomrule
\end{tabular}
\end{center}

\begin{quote}
换句话说：\textbf{只要不加限幅，这个反解在数值上就是恒等式}。
真正让"精确线性化"失效的不是数学，而是执行器——这一点在后续每一讲都会反复出现。
\end{quote}

\section{奇异点：摆杆通道会失去权限}

式 \eqref{eq:after} 中 $\cos\theta$ 是 $v$ 的系数。当 $\theta\to\pm90^\circ$ 时：

\begin{center}
\small
\begin{tabular}{@{}ccc@{}}
\toprule
$\theta$ & $\cos\theta$ & 小车加速度对 $\ddot\theta$ 的影响 \\
\midrule
$0^\circ$ & $1.00$ & 满权限 \\
$30^\circ$ & $0.87$ & $87\%$ \\
$60^\circ$ & $0.50$ & 一半 \\
$80^\circ$ & $0.17$ & 几乎失效 \\
$90^\circ$ & $0$ & \textbf{完全失效} \\
\bottomrule
\end{tabular}
\end{center}

这个表解释了一件重要的事：\textbf{摆杆水平时，无论小车怎么加速，都改变不了它的角加速度}。
所以在第 5 讲的滑模控制器里，我们对 $\cos\theta$ 做了钳位（$|\cos\theta|\ge0.25$），
并在 $\theta$ 较大时根本不使用依赖 $1/\cos\theta$ 的通道——那段交给摆起控制器。

\section{"换坐标"到底带来多少好处}

一个直接对照实验：同样的能量整定律 $v = k_E(E-E_d)\,\mathrm{sat}(\dot\theta\cos\theta/\phi)$，

\begin{center}
\small
\begin{tabular}{@{}lccc@{}}
\toprule
实现方式 & 最佳 $k_E$ & 能把摆送到顶的比例 & 说明 \\
\midrule
直接调力（无力矩补偿） & $6.0$ & $0/3$ & 冷启动时摩擦与耦合项吃掉能量 \\
加速度域（式 \eqref{eq:inv}） & $6.0$ & $3/3$ & 同一增益下全部成功 \\
\bottomrule
\end{tabular}
\end{center}

同一个控制律、同一个增益，只因为中间加了一层精确反解，就从"一个都上不去"变成"全部上去"。
这不是玄学：能量泵的推导假设"$v$ 可以任意指定"，而直接调力时这个假设根本不成立。

\keypoint{这就是"反馈线性化"的实际价值：它不是为了炫技，
而是为了\textbf{让外环控制律的推导前提成立}。}

\section{本讲小结}

\begin{enumerate}[leftmargin=2em,itemsep=2pt]
  \item 小车通道完全驱动 $\Rightarrow$ 式 \eqref{eq:inv} 是精确反解，全域可逆。
  \item 反解后系统变成 $p^{(2)}=v$ 加一个仍非线性的摆杆方程 \eqref{eq:after}。
  \item 精度验证：不作限幅时反解误差 $\sim10^{-4}$，本质上是恒等式。
  \item 奇异点：$\cos\theta\to0$ 时摆杆通道失去权限，必须显式处理。
  \item 同一能量整定律，加了反解才能工作（$0/3\to3/3$）。
\end{enumerate}

\section*{习题}
\addcontentsline{toc}{section}{习题}
\begin{enumerate}[leftmargin=2em,itemsep=6pt]
  \item 从式 \eqref{eq:pdd_raw} 出发，独立推导 \eqref{eq:inv}，并验证各项符号。
  \item 若改为指定摆杆加速度（令 $\ddot\theta=w$），反解会多出一个什么奇异项？
        在什么角度附近数值上会爆炸？
  \item 在式 \eqref{eq:inv} 中逐项去掉"干扰"（摩擦、科氏、重力耦合），
        各做一次仿真，看能量泵的成功率下降多少。
  \item 把 $\cos\theta$ 的钳位从 $0.25$ 改成 $0.1$ 与 $0.5$，观察滑模控制器给不出指令的角度范围。
  \item （编程）用 $10^{-2}$ 到 $10^{-5}$ 的步长验证反解精度，画误差随步长的收敛曲线。
\end{enumerate}

\section*{延伸阅读}
\addcontentsline{toc}{section}{延伸阅读}
\begin{itemize}[leftmargin=2em,itemsep=2pt]
  \item J.-J. Slotine, W. Li, \emph{Applied Nonlinear Control}, 第 6 章（输入–输出线性化与零动态）。
  \item S. Fantoni, R. Lozano, \emph{Non-linear Control for Underactuated Mechanical Systems}, 第 2 章。
\end{itemize}
